\documentclass[../../main.tex]{subfiles}
\begin{document}

\begin{center}
  {\fontsize{48pt}{58pt}\selectfont \textbf{1981}\normalsize\textbf{年度}}
\end{center}

\begin{table}[htb]
  \centering
  \caption{出題テーマと難易度}
  \begin{tabular}{|c|p{6.2cm}|c|c|c|}
    \hline
    \textbf{問題番号} & \textbf{テーマ} & \textbf{思考量} & \textbf{計算量} & \textbf{総合} \\
    \hline
    第1問 & 互いに交わらない二元集合の数 & 難 & 普 & 難 \\
    \hline
    第2問 & 正六角形の頂点を選ぶ三角形の期待面積 & 普 & 普 & 普 \\
    \hline
    第3問 & 放物線の法線上の点の軌跡と面積 & 普 & 普 & 普 \\
    \hline
    第4問 & 空間内の角度条件を満たす点の個数 & 難 & 難 & 難 \\
    \hline
    第5問 & 正多角錐の体積の最大値と極限 & 普 & 難 & 普 \\
    \hline
    第6問 & 三次関数の不等式条件と積分の最小化 & 普 &  & 難 \\
    \hline
  \end{tabular}
\end{table}

\begin{table}[htb]
  \centering
  \caption{解答}
  \renewcommand{\arraystretch}{1.4}
  \begin{tabular}{|c|c|p{9cm}|}
    \hline
    \textbf{問題番号} & \textbf{小問} & \textbf{答え} \\
    \hline
    第1問 & -- & $n\ge2k$なら$f(n,k)=\dfrac{n!}{2^kk!(n-2k)!}$，$n<2k$なら$0$．$f(n,k)=f(n,1)$となるのは$(n,k)=(6,3)$ \\
    \hline
    第2問 & -- & 期待値は$\dfrac{\sqrt{3}}{4}$ \\
    \hline
    第3問 & -- & 面積は$\dfrac{1}{3}$ \\
    \hline
    第4問 & -- & $t=\tan\alpha$としたとき，$t>\sqrt{2}$では$0$個，$\alpha=0$，$\alpha=\dfrac{\pi}{4}$，$t=\sqrt{2}$では$1$個，その他の$0<t<\sqrt{2}$では$2$個（各$z$座標は解答参照） \\
    \hline
    第5問 & (1) & $v_n=\dfrac{4\sqrt{5}}{375}n\tan\dfrac{\pi}{n}$ \\
    \cline{2-3}
    & (2) & 極限値は$\dfrac{4\sqrt{5}}{375}\pi$ \\
    \hline
    第6問 & (1) & $c=-a$，$d=-b$，$b\le\min\left\{-1,-|a|-\dfrac{1}{2}\right\}$ \\
    \cline{2-3}
    & (2) & $f(x)=1-x^2$ \\
    \hline
  \end{tabular}
\end{table}

\paragraph{試験形態}
6問．試験時間は要確認．

\end{document}
